\section{Fondements de la relativité restreinte}

    \subsection{Postulats d’Einstein}

    \begin{omed}{Postulats}
        \begin{enumerate}
            \item Toutes les lois de la physique sont les mêmes dans tout référentiel inertiel.
            \item Dans le vide, la lumière se propage toujours avec la même vitesse, indépendamment de la vitesse de la source.
        \end{enumerate}
    \end{omed}

    Considérons deux référentiels inertiels $\mathcal{R}$ et $\mathcal{R}'$. Un événement $\mathcal{E}$ est repéré par ses coordonnées $(t,x,y,z)$ dans $\mathcal{R}$ et $(t', x', y', z')$ dans $\mathcal{R}'$. Pour simplifier, on considèrera que les origines spatiales des deux référentiels $O$ et $O'$ coïncident, ainsi que leurs origines temporelles $t = 0$ et $t' = 0$.

    Considérons un flash de lumière émis en $t = 0$ à l’origine $O$ de $\mathcal{R}$. Un signal lumineux se propage sphériquement autour de l’origine à la vitesse de la lumière $c$ :
    \[ c^2 t^2 - x^2 - y^2 - ^2 = 0 \]   
    Dans le référentiel $\mathcal{R}'$, la même équation est vérifiée :
    \[ c^2 t'^2 - x'^2 - y'^2 - z'^2 = 0 \]
    En remplaçant les coordonnées de $\mathcal{R}'$ par leurs expressions dans $\mathcal{R}$, on définit la forme quadratique :
    \begin{align*}
        Q(t,x,y,z) 
        &:= c^2 t'^2 - x'^2 - y'^2 - z'^2 \\
        &= \sum_{i,j} \mathcal{A}_{ij} x^i x^j + \sum_{i} \mathcal{B}_{i} x^i(ct) + \mathcal{C} (ct)^2 \\
    \end{align*}
    où les coefficients $\mathcal{A}_{ij}$, $\mathcal{B}_i$ et $\mathcal{C}$ dépendent des référentiels $\mathcal{R}$ et $\mathcal{R}'$ et où les $x^i$ représentent les trois coordonnées spaciales.

    Tout point spacial de coordonneés $(x,y,z)$ sera atteint par le signal lumineux à l’intant $t_{\ell}$ défini par $ct_{\ell} = \sqrt{x^2 + y^2 + z^2}$, valeur pour laquelle la forme quadratique $Q(t,x,y,z)$ s’annule. Le signal atteignant au même instant les points $(x,y,z)$ et $(-x,-y,-z)$ symétriques par rapport à $O$, on sait aussi que $Q(t_{\ell}, -x, -y, -z) = 0$, ce qui montre que les coefficients $\mathcal{B}_i$ sont nuls.

\begin{comment}
    Si on suppose de plus que l’espace est \textit{isotrope}, on remarque que la matrice $A = (\mathcal{A}_{ij})$ doit être proportionnelle à l’identité (pas de direction privilégiée). Ainsi, $Q$ s’exprime nécessairement sous la forme 
    \[ Q(t,x,y,z) = \mathcal{C}\left((ct)^2 - x^2 - y^2 - z^2\right) \]   
    Sachant que la seule grandeur qui caractérise ces deux référentiels \textit{relativement} l’un à l’autre est la vitesse relative $\mathbf{v}$, et par isotropie de l’espace, on voit que $\mathcal{C}$ ne peut dépendre que de la norme de cette vitesse relative. Pour obtenir sa valeur, on remarque que 
    \[ \mathcal{C}\left(\nor{\mathbf{v}_{\mathcal{R}''/\mathcal{R}}}\right) = \mathcal{C}\left(\nor{\mathbf{v}_{\mathcal{R}''/\mathcal{R}'}}\right)\mathcal{C}\left(\nor{\mathbf{v}_{\mathcal{R}'/\mathcal{R}}}\right) \]
    La norme de $\mathbf{v}_{\mathcal{R}''/\mathcal{R}}$ dépendant non seulement des deux autres normes mais aussi de l’angle entre celles-ci, on a nécessairement $\mathcal{C} = 1$. Les deux postulats d’Einstein, combinés aux hypothèses d’homogénéité et d’isotropie de l’espace, conduisent donc à la relation
    \[ c^2 t^2 - x^2 - y^2 - z^2 = c^2 t'^2 - x'^2 - y'^2 - z'^2 \]
    entre les coordonnées de tout événement repéré dans $\mathcal{R}$ et $\mathcal{R}'$. Ces tranformations de coordonnées sont appelées \textbf{transformations de Lorentz}.
\end{comment}

    \subsection{Le groupe Lorentz}

    \begin{omed}{Définition (transformation de Lorentz)}
        On appelle \textbf{transformation de Lorentz} toute transformation linéaire 
        \begin{align*}
            x'^{\mu} &= \sum_{\nu = 0}^3 \Lambda_{\nu}^{\mu} x^{\nu} \\
            x'^{\mu} &\overset{\Updownarrow}{=} \symbf{\Lambda} x^{\nu} 
        \end{align*}
        reliant par la matrice carrée $\symbf{\Lambda}$ les quadrivecteurs $x'^{\mu} = (ct', x', y', z')$ et $x^{\mu} = (ct,x,y,z)$ vérifiant la relation 
        \begin{align*}
            c^2 t^2 - x^2 - y^2 - z^2 &= c^2 t'^2 - x'^2 - y'^2 - z'^2 \\
            x'^{\mu\top} \mathbf{G} x'^{\mu} &\overset{\Updownarrow}{=} x^{\mu\top} \mathbf{G} x^{\mu} 
        \end{align*}
        où $\mathbf{G} = \text{diag}(1, -1, -1, -1)$.
    \end{omed}

\begin{comment}
    \begin{omed}{Proposition (groupe de Lorentz)}
        Les matrices $\symbf{\Lambda}$ satisfaisant les conditions ci-dessus forment un groupe, dit \textbf{groupe de Lorentz}.
    \end{omed}
\end{comment}

    \begin{demo}{Preuve}
        En combinant les deux égalités de la transformation de Lorentz, on obtient la relation 
        \[ x^{\mu\top} \symbf{\Lambda}^{\top} \mathbf{G} \symbf{\Lambda} x^{\mu} =  x^{\mu\top} \mathbf{G} x^{\mu} \]
        Celle-ci étant valable pour tout $x^{\mu}$, on en déduit que $ \symbf{\Lambda}^{\top} \mathbf{G}  \symbf{\Lambda} = \mathbf{G}$, ce qui permet de vérifier que ces matrices forment un groupe. Leur symétrique s’écrit $\symbf{\Lambda}^{-1} = \mathbf{G} \symbf{\Lambda}^{\top} \mathbf{G}$ (étant donné aussi que $\mathbf{G}^2 = \mathbf{I}$), et l’élément neutre est $\mathbf{I}$.
    \end{demo}

    \subsection{Transformations de Lorentz spéciales}

        \subsubsection{Transformation de Lorentz le long de l’axe des $x$}

        \begin{omed}{Transformation de Lorentz spéciale le long de l’axe $x$}
            \begin{align*}
                ct' &= \gamma(ct - \beta x) \\
                x' &= \gamma(x - \beta ct) \\
                y' &= y \\
                z' &= z \\
            \end{align*}
            où on a les coefficients adimensionnes
            \begin{align*}
                \beta &= \frac{v}{c} \\
                \gamma &= \frac{1}{\sqrt{1 - \frac{v^2}{c^2}}}
            \end{align*}
        \end{omed}

        \begin{demo}{Démonstration}
            Considérons que deux des trois coordonnées spatiales restent inchangées : $y' = y$ et $z' = z$. On a alors 
            \[ c^2 t^2 - x^2 = c^2 t'^2 - x'^2 \]   
            La transformation linéaire la plus générale reliant $(ct, x)$ à $(ct', x')$ s’écrit 
        \begin{align*}
            ct' &= \lambda_1 ct + \lambda_2 x \\
            x' &= \lambda_3 x + \lambda_4 ct \\
        \end{align*}
        Pour satisfaire la transformation de Lorentz, il faut alors que
        \begin{align*}
            \lambda_1^2 - \lambda_4^2 &= 1 \\
            \lambda_3^2 - \lambda_2^2 &= 1 \\
            \lambda_1 \lambda_2 - \lambda_3 \lambda_4 &= 0 \\
        \end{align*}
        La trajectoire de l’origine $O'$ de $\mathcal{R}'$ est caractérisée par $x' = 0$, \textit{i.e.} $\lambda_3 x + \lambda_4 ct = 0$, qui nous dit que la vitesse de $O'$ selon $\mathcal{R}$ est parallèle à l’axe des $x$ et de valeur $v = -c \frac{\lambda_4}{\lambda_3}$. La solution de ces équations est
        \begin{align*}
            \lambda_1 = \lambda_3 &= \left(1 - \frac{v^2}{c^2}\right)^{-\frac{1}{2}} := \gamma \\
            \lambda_2 = \lambda_4 &= - \frac{v}{c} \lambda_1 := - \beta \gamma \\
        \end{align*}
        \end{demo}

        Il est parfois utile d’exprimer les coefficients $\beta$ et $\gamma$ sous la forme 
        \begin{align*}
            \gamma &= \text{ch}(\alpha) \\
            \beta &= \text{th}(\alpha)
        \end{align*}
        où le paramètre $\alpha$ est appelé \textbf{rapidité}. La transformation de Lorentz s’écrit alors 
        \[ \begin{bmatrix}
            ct' \\
            x'
        \end{bmatrix} = \begin{bmatrix}
            \text{ch}(\alpha) & -\text{sh}(\alpha) \\
            -\text{sh}(\alpha) & \text{ch}(\alpha)
        \end{bmatrix} \begin{bmatrix}
            ct \\
            x
        \end{bmatrix} \]

        \subsubsection{Transformation de Lorentz d’axe quelconque}

        \begin{omed}{Transformation de Lorentz d’axe quelconque}
            Sous forme matricielle, la transformation de Lorentz spéciale de paramètre $\symbf{\beta} = \frac{\mathbf{v}}{c}$ s’écrit 
            \[ \symbf{\Lambda}_{\text{boost}}(\symbf{\beta}) = \begin{bmatrix}
                \gamma & -\gamma \beta_x & - \gamma \beta_y & - \gamma \beta_z \\
                -\gamma \beta_x & 1 + \frac{\gamma - 1}{\beta^2} \beta_x^2  & \frac{\gamma - 1}{\beta^2} \beta_x \beta_y & \frac{\gamma - 1}{\beta^2} \beta_x \beta_z \\
                -\gamma \beta_y & \frac{\gamma - 1}{\beta^2} \beta_x \beta_y & 1 + \frac{\gamma - 1}{\beta^2} \beta_y^2 & \frac{\gamma - 1}{\beta^2} \beta_y \beta_z \\
                -\gamma \beta_z & \frac{\gamma - 1}{\beta^2} \beta_x \beta_z & \frac{\gamma - 1}{\beta^2} \beta_y \beta_z & 1 + \frac{\gamma - 1}{\beta^2} \beta_z^2 \\
            \end{bmatrix} \]
        \end{omed}

    \subsection{Conséquences de la transformation de Lorentz}

        \subsubsection{Relativité de la simultanéité}

        Contrairement à la physique newtonienne, en relativité, deux événements se produisant simultanément dans un référentiel $\mathcal{R}$ ne sont plus nécessairement simultanées dans $\mathcal{R}'$. En effet, si on considère les deux événements $\mathcal{E}_1$ et $\mathcal{E}_2$ se réalisant à $t_1 = t_2$ dans $\mathcal{R}$, alors dans $\mathcal{R}'$ (dont la transformation de Lorentz est le long de l’axe $x$), 
        \begin{align*}
            \Delta t' &= t_2' - t_1' = \gamma\left(\Delta t - \frac{\beta}{c} \Delta x\right) \\
            &= - \gamma \frac{\beta}{c} \Delta x \\
        \end{align*}

        \subsubsection{Dilatation du temps}

        On suppose désormais que $\Delta x = 0$ mais que $\Delta t \neq 0$ pour deux événements $\mathcal{E}_1$ et $\mathcal{E}_2$ dans $\mathcal{R}$. On obtient cette fois-ci dans $\mathcal{R}'$ transformé selon Lorentz le long de l’axe $x$, 
        \begin{align*}
            \Delta t' 
            &= t_2' - t_1' = \gamma\left(\Delta t - \frac{\beta}{c} \Delta x\right) \\
            &= \gamma \Delta t \\
        \end{align*}
        Cet intervalle de temps est donc toujours supérieur dans $\mathcal{R}'$ : il y a \textbf{dilatation du temps}.

        \subsubsection{Contraction des longueurs}

        Intéressons-nous désormais à la longueur d’un objet : par exemple une règle immobile dans le référentiel $\mathcal{R}$ de l’observateur, et choisissons l’axe des $x$ parallèle à cette règle. Sa longueur dans $\mathcal{R}$ est définie par $L = \Delta x = x_2 - x_1$, et celle dans $\mathcal{R}'$ est $L' = \Delta x' = x'_2(t') - x'_1(t')$. Les deux extrémités dans $\mathcal{R}'$ ayant la même vitesse, $L'$ est indépendant du temps. 

        Les événements $\mathcal{E}_1$ et $\mathcal{E}_2$ simultanés dans $\mathcal{R}'$ correspondant aux positions des extrémités de la règle à l’instant $t'$ et repérés respectivement par leurs coordonnées $(t', x_1')$ et $(t', x_2')$ ne sont plus simultanés dans $\mathcal{R}$, et 
        \[ \Delta t = \beta \frac{L}{c} \]   
        On obtient alors 
        \begin{align*}
            L' 
            &= \gamma (L - \beta c \Delta t) = \gamma (1 - \beta^2) L \\
            &= \sqrt{1- \beta^2} L = \frac{L}{\gamma}
        \end{align*}
        La longueur mesurée dans un référentiel quelconque est donc toujours inférieure à celle mesurée dans le référentiel particulier où l’objet est immobile : c’est la \textbf{contraction des longueurs}.

    \subsection{Loi de composition des vitesses}

        \subsubsection{Composition de deux transformations de Lorentz de même axe}

        En notant $\symbf{\Lambda}(\alpha) = \begin{bmatrix}
            \text{ch}(\alpha) & -\text{sh}(\alpha) \\
            - \text{sh}(\alpha) & \text{ch}(\alpha) \\
        \end{bmatrix}$, et en associant à la transformation de $\mathcal{R}$ à $\mathcal{R}'$ l’angle $\alpha_1$ et à la transformation $\mathcal{R}'$ à $\mathcal{R}''$ l’angle $\alpha_2$, on constate que l’on peut écrire la transformation de Lorentz selon l’axe $x$ entre $\mathcal{R}$ et $\mathcal{R}''$ par l’équation 
        \[ \begin{bmatrix}
            ct'' \\
            x''
        \end{bmatrix} = \symbf{\Lambda}(\alpha_2)\symbf{\Lambda}(\alpha_1) \begin{bmatrix}
            ct \\
            x
        \end{bmatrix} \]   
        Or $\symbf{\Lambda}(\alpha_2)\symbf{\Lambda}(\alpha_1) = \symbf{\Lambda}(\alpha_1 + \alpha_2)$. Cela correspond donc au total à une transformation de Lorentz spéciale le long de l’axe des $x$, de vitesse $V = c \text{th}(\alpha_1 + \alpha_2) = c \frac{\text{th}(\alpha_1) + \text{th}(\alpha_2)}{1 + \text{th}(\alpha_1)\text{th}(\alpha_2)}$. La vitesse relative entre $\mathcal{R}$ et $\mathcal{R}'$ est donc donnée par 
        \[ v_{\mathcal{R}''/\mathcal{R}} = \frac{v_{\mathcal{R}''/\mathcal{R}'} + v_{\mathcal{R}'/\mathcal{R}}}{1 + \frac{v_{\mathcal{R}''/\mathcal{R}'}v_{\mathcal{R}'/\mathcal{R}}}{c^2}} \]
        Dans le cas de vitesse faibles devant celle de la lumière, on retrouve la loi d’addition des vitesses de la physique newtonnienne.

        \subsubsection{Composition des vitesses}

        On considère maintenant le cas d’une particule animée d’un mouvement quelconque, de vitesse $\mathbf{w}$ et $\mathbf{w}'$ dans les référentiels intertiels $\mathcal{R}$ et $\mathcal{R}'$. En utilisant la transformation de Lorentz entre $\mathcal{R}$ et $\mathcal{R}'$, on obtient 
        \begin{align*}
            w_x' 
            &= \frac{\text{d}x'}{\text{d}t'} = \frac{\gamma(dx - \beta c dt)}{\gamma(cdt - \beta dx)} \\
            &= \frac{w_x - v}{1 - \frac{v w_x}{c^2}} \\
        \end{align*}
        De même, on a pour les composantes orthogonales à l’axe de la tranformation 
        \begin{align*}
            w_y' &= \frac{\text{d}y'}{\text{d}t'} = \frac{w_y}{\gamma\left(1 - \frac{\beta}{c} w_x\right)} \\
            w_z' &= \frac{\text{d}z'}{\text{d}t'} = \frac{w_z}{\gamma\left(1 - \frac{\beta}{c} w_x\right)} \\
        \end{align*}

        Pour une transformation d’axe quelconque, on peut généraliser en 
        \begin{align*}
            \mathbf{w}_{||}' &= \frac{\mathbf{w}_{||} - \mathbf{v}}{1 - \frac{\mathbf{v} \cdot \mathbf{w}}{c^2}} \\
            \mathbf{w}_{\perp}' &= \frac{\mathbf{w}_{\perp}}{\gamma(1 - \frac{\mathbf{v} \cdot \mathbf{w}}{c^2})} \\
        \end{align*}


    

        